% ============================== CHƯƠNG 3 ==============================
\chapter{Dữ liệu và tiền xử lý}
\label{chap:data}

\section{Mô tả bộ dữ liệu}
\todo{Nguồn Kaggle DataCo~\cite{dataco2019}, 53 cột, nhóm thuộc tính; bảng thống kê mô tả.}

\section{Khảo sát chất lượng dữ liệu (Bronze profiling)}
\begin{table}[H]
\centering
\caption{Các vấn đề chất lượng dữ liệu và cách xử lý}
\label{tab:dq-issues}
\begin{tabularx}{\linewidth}{lXX}
\toprule
\textbf{\#} & \textbf{Vấn đề} & \textbf{Giải pháp} \\ \midrule
1 & 6 địa danh lỗi font Latin-1 & Từ điển \texttt{CITY\_REPAIR\_MAP} \\
2 & 5 cặp cột trùng lặp 100\% & Loại bỏ cột dư thừa \\
3 & \texttt{Product Description} 100\% null, \texttt{Order Zipcode} 86,24\% null & Loại bỏ \\
4 & 6 cột PII (email, mật khẩu, họ tên, địa chỉ, ảnh) & Loại bỏ trước khi lên Cloud \\
5 & Ngày dạng chuỗi \texttt{M/D/YYYY H:MM} & Ép kiểu timestamp \\
6 & Thiếu biến nghiệp vụ & Feature engineering 9 cột \\
\bottomrule
\end{tabularx}
\end{table}

\section{Làm sạch và làm giàu dữ liệu (Silver)}
\todo{Mô tả các bước trong \texttt{etl/step02\_transform\_silver.py}; bảng 9 feature mới.}

\section{Kiểm định chất lượng dữ liệu tự động}
Quality gate gồm 43 kiểm tra (row count, schema, null, FK/PK, business rules, outlier, temporal):
41 PASS, 0 FAIL, 2 WARN.
\todo{Bảng kết quả từ \texttt{data/gold/data\_quality\_report.json}.}
